%% ---------------------------------------------------------------------
\section{Señales de prueba}
\label{sec:er-senales}
%% ---------------------------------------------------------------------

El \clave{error en estado estacionario} es el límite del error de seguimiento cuando $t\to\infty$, ante una señal de prueba. Las tres tradicionales son (\cref{fig:er-senales}):
\begin{itemize}
  \item \textbf{Escalón}, $r(t)=\escalon(t)$, $R(s)=1/s$: mantener una posición o un valor constante (un satélite geoestacionario, una temperatura).
  \item \textbf{Rampa}, $r(t)=t\,\escalon(t)$, $R(s)=1/s^2$: seguir un blanco que se mueve a velocidad constante (una antena que sigue un satélite).
  \item \textbf{Parábola}, $r(t)=\frac{t^2}{2}\escalon(t)$, $R(s)=1/s^3$: seguir un blanco con aceleración constante.
\end{itemize}

\begin{figure}[ht]
\centering
\includegraphics[width=\linewidth]{cap06/f_senales}
\caption{Las tres señales de prueba.}
\label{fig:er-senales}
\end{figure}

\begin{cuidado}[Primero la estabilidad]
El teorema del valor final solo vale si $sE(s)$ tiene todos sus polos en el semiplano izquierdo, es decir, si el lazo cerrado es estable. Calcular un ``error en estado estacionario'' de un lazo inestable da un número que no significa nada. \alerta{Verifique la estabilidad primero} (\cref{cap:es}).
\end{cuidado}
